% ---------------------------------------------------------------------------
% Prezentace k tématu 01 – HW platformy a vývojové prostředí
% Část 1 ze 2: učivo U1 a U2 (+ zavedení U5, U6, U7)
%
% ŠVP: V1, V3 | Učivo: U1 (možnosti tvorby programů pro různé HW platformy),
% U2 (vybrané platformy AVR a ARM, přednosti a omezení),
% okrajově U5 (kompilace), U6 (bootloader, ICSP), U7 (rozdělení paměti).
%
% Přehledové téma – jde o orientaci v krajině, ne o detaily. Každý slide nese
% jednu myšlenku; radši víc krátkých slidů než jeden nabitý.
%
% ---------------------------------------------------------------------------
% KLÍČOVÉ SNÍMKY
% ---------------------------------------------------------------------------
% Prostředí `keyframe` (červený pruh vlevo) je vyhrazené pro snímky, které
% drží strukturu hodiny. V této prezentaci je jich pět:
%   slide  2  Co dnes zvládneme            – cíle hodiny
%   slide  5  Část 1 (U1): co se dozvíte   – cíle první části
%   slide 20  Shrnutí části 1 (U1)
%   slide 21  Část 2 (U2): co se dozvíte   – cíle druhé části
%   slide 33  Shrnutí hodiny               – závěr + co příště
% Dělicí snímky sekcí jsou proto vypnuté (\SPSSectionSlidesOff); \section
% zůstává jen pro záložky v PDF.
%
% Slide „Čipy zblízka`` je REZERVA – když se nestíhá, vypustit.
%
% Časování (45 min):
%   0-3    slidy 1-2    cíle hodiny
%   3-6    slidy 3-4    kde všude je mikrokontroler, proč se tím zabýváme
%   6-8    slide 5      cíle části U1
%   8-26   slidy 6-19   U1: platforma, instrukční sada, cesta kódu, úrovně
%                       programování, HAL, paměti, nahrávání
%   26-28  slide 20     shrnutí U1
%   28-30  slide 21     cíle části U2
%   30-42  slidy 22-32  U2: desky, AVR, ARM, architektura a registry,
%                       srovnání, kvíz, časté záměny
%   42-45  slide 33     shrnutí hodiny
%
% Snímky 3 a 22 jsou celostránkové fotografie. Šířka 0.82\textwidth je
% maximum, které se vejde pod titulek; při \textwidth se obrázek dole ořízne
% (beamer to nehlásí).
%
% Překlad:  latexmk -xelatex prezentace.tex
% ---------------------------------------------------------------------------
\documentclass[12pt,aspectratio=169]{beamer}

% ---------------------------------------------------------------------------
% Společná preambule pro prezentace předmětu Programování (3. ročník)
%
% Vzhled řeší motiv beamerthemeSPSProsek.sty ve stejné složce (převzatý
% z šablony předmětů ZPG/SNT na FS ČVUT a přeznačkovaný na SPŠ na Proseku).
% Tady se jen načte a doplní se metadata a makra předmětu.
%
% Překlad: XeLaTeX (kvůli fontspec a českým znakům)
%   latexmk -xelatex prezentace.tex
% nebo jednoduše  .\nastroje\build_prezentace.ps1
%
% Vkládá se z prezentace v temata/<tema>/prezentace/ takto:
%   % ---------------------------------------------------------------------------
% Společná preambule pro prezentace předmětu Programování (3. ročník)
%
% Vzhled řeší motiv beamerthemeSPSProsek.sty ve stejné složce (převzatý
% z šablony předmětů ZPG/SNT na FS ČVUT a přeznačkovaný na SPŠ na Proseku).
% Tady se jen načte a doplní se metadata a makra předmětu.
%
% Překlad: XeLaTeX (kvůli fontspec a českým znakům)
%   latexmk -xelatex prezentace.tex
% nebo jednoduše  .\nastroje\build_prezentace.ps1
%
% Vkládá se z prezentace v temata/<tema>/prezentace/ takto:
%   % ---------------------------------------------------------------------------
% Společná preambule pro prezentace předmětu Programování (3. ročník)
%
% Vzhled řeší motiv beamerthemeSPSProsek.sty ve stejné složce (převzatý
% z šablony předmětů ZPG/SNT na FS ČVUT a přeznačkovaný na SPŠ na Proseku).
% Tady se jen načte a doplní se metadata a makra předmětu.
%
% Překlad: XeLaTeX (kvůli fontspec a českým znakům)
%   latexmk -xelatex prezentace.tex
% nebo jednoduše  .\nastroje\build_prezentace.ps1
%
% Vkládá se z prezentace v temata/<tema>/prezentace/ takto:
%   \input{../../../spolecne/styl/preambule}
% ---------------------------------------------------------------------------

\usepackage{iftex}

% --- písma a čeština -------------------------------------------------------
\ifXeTeX
  \usepackage{fontspec}
  \usepackage{polyglossia}
  \setmainlanguage{czech}
  \setmainfont{Latin Modern Roman}
  \setsansfont{Latin Modern Sans}
  \setmonofont{Latin Modern Mono}[Scale=MatchLowercase]
\else
  % Záložní varianta, kdyby někdo přeložil pdfLaTeXem.
  \usepackage[T1]{fontenc}
  \usepackage[utf8]{inputenc}
  \usepackage{lmodern}
  \usepackage[czech]{babel}
\fi

\usepackage{graphicx}
\usepackage{booktabs}
\usepackage{xcolor}
\usepackage{tikz}
\usetikzlibrary{arrows.meta, positioning, shapes.geometric}

% --- motiv SPŠ na Proseku --------------------------------------------------
% Cesta ke složce se stylem a logy. Počítá s tím, že prezentace leží
% v temata/<tema>/prezentace/; zkusí se ale i mělčí zanoření, aby šlo
% přeložit i rozpracovaný soubor jinde v repozitáři.
\IfFileExists{../../../spolecne/styl/beamerthemeSPSProsek.sty}
  {\newcommand{\SPSStylPath}{../../../spolecne/styl}}
  {\IfFileExists{../../spolecne/styl/beamerthemeSPSProsek.sty}
    {\newcommand{\SPSStylPath}{../../spolecne/styl}}
    {\IfFileExists{../spolecne/styl/beamerthemeSPSProsek.sty}
      {\newcommand{\SPSStylPath}{../spolecne/styl}}
      {\newcommand{\SPSStylPath}{.}}}}
\usepackage{\SPSStylPath/beamerthemeSPSProsek}

\setbeamertemplate{footline}[SPSProsek]
\beamertemplatenavigationsymbolsempty

% --- velikost písma ve výpisech kódu ---------------------------------------
% Výchozí je \footnotesize. Když se dlouhá ukázka nevejde na slide, zmenši
% ji buď globálně tady, nebo jen v jedné prezentaci stejným řádkem za
% \input{...preambule}:
% \renewcommand{\SPSCodeSize}{\scriptsize}

% --- zkratky pro vkládání kódu ---------------------------------------------
\newcommand{\kod}[1]{\texttt{\small #1}}          % kód uvnitř věty
\newcommand{\kodSoubor}[1]{\lstinputlisting{#1}}  % celá skica ze souboru
\newcommand{\kodSouborRadky}[3]{\lstinputlisting[firstline=#2,lastline=#3]{#1}}

% --- upozornění na slidu ---------------------------------------------------
\newcommand{\pozor}[1]{\textcolor{akcent}{\textbf{#1}}}

% --- metadata předmětu -----------------------------------------------------
\author{Jaroslav Bušek}
\institute{Střední průmyslová škola na Proseku}
% \date{} zůstává prázdné záměrně -- datum a čas posledního překladu se
% zobrazuje automaticky v patě každého slidu (\SPSVersionStamp), takže je
% při promítání hned vidět, jak stará je verze PDF.
\date{}

% ---------------------------------------------------------------------------

\usepackage{iftex}

% --- písma a čeština -------------------------------------------------------
\ifXeTeX
  \usepackage{fontspec}
  \usepackage{polyglossia}
  \setmainlanguage{czech}
  \setmainfont{Latin Modern Roman}
  \setsansfont{Latin Modern Sans}
  \setmonofont{Latin Modern Mono}[Scale=MatchLowercase]
\else
  % Záložní varianta, kdyby někdo přeložil pdfLaTeXem.
  \usepackage[T1]{fontenc}
  \usepackage[utf8]{inputenc}
  \usepackage{lmodern}
  \usepackage[czech]{babel}
\fi

\usepackage{graphicx}
\usepackage{booktabs}
\usepackage{xcolor}
\usepackage{tikz}
\usetikzlibrary{arrows.meta, positioning, shapes.geometric}

% --- motiv SPŠ na Proseku --------------------------------------------------
% Cesta ke složce se stylem a logy. Počítá s tím, že prezentace leží
% v temata/<tema>/prezentace/; zkusí se ale i mělčí zanoření, aby šlo
% přeložit i rozpracovaný soubor jinde v repozitáři.
\IfFileExists{../../../spolecne/styl/beamerthemeSPSProsek.sty}
  {\newcommand{\SPSStylPath}{../../../spolecne/styl}}
  {\IfFileExists{../../spolecne/styl/beamerthemeSPSProsek.sty}
    {\newcommand{\SPSStylPath}{../../spolecne/styl}}
    {\IfFileExists{../spolecne/styl/beamerthemeSPSProsek.sty}
      {\newcommand{\SPSStylPath}{../spolecne/styl}}
      {\newcommand{\SPSStylPath}{.}}}}
\usepackage{\SPSStylPath/beamerthemeSPSProsek}

\setbeamertemplate{footline}[SPSProsek]
\beamertemplatenavigationsymbolsempty

% --- velikost písma ve výpisech kódu ---------------------------------------
% Výchozí je \footnotesize. Když se dlouhá ukázka nevejde na slide, zmenši
% ji buď globálně tady, nebo jen v jedné prezentaci stejným řádkem za
% % ---------------------------------------------------------------------------
% Společná preambule pro prezentace předmětu Programování (3. ročník)
%
% Vzhled řeší motiv beamerthemeSPSProsek.sty ve stejné složce (převzatý
% z šablony předmětů ZPG/SNT na FS ČVUT a přeznačkovaný na SPŠ na Proseku).
% Tady se jen načte a doplní se metadata a makra předmětu.
%
% Překlad: XeLaTeX (kvůli fontspec a českým znakům)
%   latexmk -xelatex prezentace.tex
% nebo jednoduše  .\nastroje\build_prezentace.ps1
%
% Vkládá se z prezentace v temata/<tema>/prezentace/ takto:
%   \input{../../../spolecne/styl/preambule}
% ---------------------------------------------------------------------------

\usepackage{iftex}

% --- písma a čeština -------------------------------------------------------
\ifXeTeX
  \usepackage{fontspec}
  \usepackage{polyglossia}
  \setmainlanguage{czech}
  \setmainfont{Latin Modern Roman}
  \setsansfont{Latin Modern Sans}
  \setmonofont{Latin Modern Mono}[Scale=MatchLowercase]
\else
  % Záložní varianta, kdyby někdo přeložil pdfLaTeXem.
  \usepackage[T1]{fontenc}
  \usepackage[utf8]{inputenc}
  \usepackage{lmodern}
  \usepackage[czech]{babel}
\fi

\usepackage{graphicx}
\usepackage{booktabs}
\usepackage{xcolor}
\usepackage{tikz}
\usetikzlibrary{arrows.meta, positioning, shapes.geometric}

% --- motiv SPŠ na Proseku --------------------------------------------------
% Cesta ke složce se stylem a logy. Počítá s tím, že prezentace leží
% v temata/<tema>/prezentace/; zkusí se ale i mělčí zanoření, aby šlo
% přeložit i rozpracovaný soubor jinde v repozitáři.
\IfFileExists{../../../spolecne/styl/beamerthemeSPSProsek.sty}
  {\newcommand{\SPSStylPath}{../../../spolecne/styl}}
  {\IfFileExists{../../spolecne/styl/beamerthemeSPSProsek.sty}
    {\newcommand{\SPSStylPath}{../../spolecne/styl}}
    {\IfFileExists{../spolecne/styl/beamerthemeSPSProsek.sty}
      {\newcommand{\SPSStylPath}{../spolecne/styl}}
      {\newcommand{\SPSStylPath}{.}}}}
\usepackage{\SPSStylPath/beamerthemeSPSProsek}

\setbeamertemplate{footline}[SPSProsek]
\beamertemplatenavigationsymbolsempty

% --- velikost písma ve výpisech kódu ---------------------------------------
% Výchozí je \footnotesize. Když se dlouhá ukázka nevejde na slide, zmenši
% ji buď globálně tady, nebo jen v jedné prezentaci stejným řádkem za
% \input{...preambule}:
% \renewcommand{\SPSCodeSize}{\scriptsize}

% --- zkratky pro vkládání kódu ---------------------------------------------
\newcommand{\kod}[1]{\texttt{\small #1}}          % kód uvnitř věty
\newcommand{\kodSoubor}[1]{\lstinputlisting{#1}}  % celá skica ze souboru
\newcommand{\kodSouborRadky}[3]{\lstinputlisting[firstline=#2,lastline=#3]{#1}}

% --- upozornění na slidu ---------------------------------------------------
\newcommand{\pozor}[1]{\textcolor{akcent}{\textbf{#1}}}

% --- metadata předmětu -----------------------------------------------------
\author{Jaroslav Bušek}
\institute{Střední průmyslová škola na Proseku}
% \date{} zůstává prázdné záměrně -- datum a čas posledního překladu se
% zobrazuje automaticky v patě každého slidu (\SPSVersionStamp), takže je
% při promítání hned vidět, jak stará je verze PDF.
\date{}
:
% \renewcommand{\SPSCodeSize}{\scriptsize}

% --- zkratky pro vkládání kódu ---------------------------------------------
\newcommand{\kod}[1]{\texttt{\small #1}}          % kód uvnitř věty
\newcommand{\kodSoubor}[1]{\lstinputlisting{#1}}  % celá skica ze souboru
\newcommand{\kodSouborRadky}[3]{\lstinputlisting[firstline=#2,lastline=#3]{#1}}

% --- upozornění na slidu ---------------------------------------------------
\newcommand{\pozor}[1]{\textcolor{akcent}{\textbf{#1}}}

% --- metadata předmětu -----------------------------------------------------
\author{Jaroslav Bušek}
\institute{Střední průmyslová škola na Proseku}
% \date{} zůstává prázdné záměrně -- datum a čas posledního překladu se
% zobrazuje automaticky v patě každého slidu (\SPSVersionStamp), takže je
% při promítání hned vidět, jak stará je verze PDF.
\date{}

% ---------------------------------------------------------------------------

\usepackage{iftex}

% --- písma a čeština -------------------------------------------------------
\ifXeTeX
  \usepackage{fontspec}
  \usepackage{polyglossia}
  \setmainlanguage{czech}
  \setmainfont{Latin Modern Roman}
  \setsansfont{Latin Modern Sans}
  \setmonofont{Latin Modern Mono}[Scale=MatchLowercase]
\else
  % Záložní varianta, kdyby někdo přeložil pdfLaTeXem.
  \usepackage[T1]{fontenc}
  \usepackage[utf8]{inputenc}
  \usepackage{lmodern}
  \usepackage[czech]{babel}
\fi

\usepackage{graphicx}
\usepackage{booktabs}
\usepackage{xcolor}
\usepackage{tikz}
\usetikzlibrary{arrows.meta, positioning, shapes.geometric}

% --- motiv SPŠ na Proseku --------------------------------------------------
% Cesta ke složce se stylem a logy. Počítá s tím, že prezentace leží
% v temata/<tema>/prezentace/; zkusí se ale i mělčí zanoření, aby šlo
% přeložit i rozpracovaný soubor jinde v repozitáři.
\IfFileExists{../../../spolecne/styl/beamerthemeSPSProsek.sty}
  {\newcommand{\SPSStylPath}{../../../spolecne/styl}}
  {\IfFileExists{../../spolecne/styl/beamerthemeSPSProsek.sty}
    {\newcommand{\SPSStylPath}{../../spolecne/styl}}
    {\IfFileExists{../spolecne/styl/beamerthemeSPSProsek.sty}
      {\newcommand{\SPSStylPath}{../spolecne/styl}}
      {\newcommand{\SPSStylPath}{.}}}}
\usepackage{\SPSStylPath/beamerthemeSPSProsek}

\setbeamertemplate{footline}[SPSProsek]
\beamertemplatenavigationsymbolsempty

% --- velikost písma ve výpisech kódu ---------------------------------------
% Výchozí je \footnotesize. Když se dlouhá ukázka nevejde na slide, zmenši
% ji buď globálně tady, nebo jen v jedné prezentaci stejným řádkem za
% % ---------------------------------------------------------------------------
% Společná preambule pro prezentace předmětu Programování (3. ročník)
%
% Vzhled řeší motiv beamerthemeSPSProsek.sty ve stejné složce (převzatý
% z šablony předmětů ZPG/SNT na FS ČVUT a přeznačkovaný na SPŠ na Proseku).
% Tady se jen načte a doplní se metadata a makra předmětu.
%
% Překlad: XeLaTeX (kvůli fontspec a českým znakům)
%   latexmk -xelatex prezentace.tex
% nebo jednoduše  .\nastroje\build_prezentace.ps1
%
% Vkládá se z prezentace v temata/<tema>/prezentace/ takto:
%   % ---------------------------------------------------------------------------
% Společná preambule pro prezentace předmětu Programování (3. ročník)
%
% Vzhled řeší motiv beamerthemeSPSProsek.sty ve stejné složce (převzatý
% z šablony předmětů ZPG/SNT na FS ČVUT a přeznačkovaný na SPŠ na Proseku).
% Tady se jen načte a doplní se metadata a makra předmětu.
%
% Překlad: XeLaTeX (kvůli fontspec a českým znakům)
%   latexmk -xelatex prezentace.tex
% nebo jednoduše  .\nastroje\build_prezentace.ps1
%
% Vkládá se z prezentace v temata/<tema>/prezentace/ takto:
%   \input{../../../spolecne/styl/preambule}
% ---------------------------------------------------------------------------

\usepackage{iftex}

% --- písma a čeština -------------------------------------------------------
\ifXeTeX
  \usepackage{fontspec}
  \usepackage{polyglossia}
  \setmainlanguage{czech}
  \setmainfont{Latin Modern Roman}
  \setsansfont{Latin Modern Sans}
  \setmonofont{Latin Modern Mono}[Scale=MatchLowercase]
\else
  % Záložní varianta, kdyby někdo přeložil pdfLaTeXem.
  \usepackage[T1]{fontenc}
  \usepackage[utf8]{inputenc}
  \usepackage{lmodern}
  \usepackage[czech]{babel}
\fi

\usepackage{graphicx}
\usepackage{booktabs}
\usepackage{xcolor}
\usepackage{tikz}
\usetikzlibrary{arrows.meta, positioning, shapes.geometric}

% --- motiv SPŠ na Proseku --------------------------------------------------
% Cesta ke složce se stylem a logy. Počítá s tím, že prezentace leží
% v temata/<tema>/prezentace/; zkusí se ale i mělčí zanoření, aby šlo
% přeložit i rozpracovaný soubor jinde v repozitáři.
\IfFileExists{../../../spolecne/styl/beamerthemeSPSProsek.sty}
  {\newcommand{\SPSStylPath}{../../../spolecne/styl}}
  {\IfFileExists{../../spolecne/styl/beamerthemeSPSProsek.sty}
    {\newcommand{\SPSStylPath}{../../spolecne/styl}}
    {\IfFileExists{../spolecne/styl/beamerthemeSPSProsek.sty}
      {\newcommand{\SPSStylPath}{../spolecne/styl}}
      {\newcommand{\SPSStylPath}{.}}}}
\usepackage{\SPSStylPath/beamerthemeSPSProsek}

\setbeamertemplate{footline}[SPSProsek]
\beamertemplatenavigationsymbolsempty

% --- velikost písma ve výpisech kódu ---------------------------------------
% Výchozí je \footnotesize. Když se dlouhá ukázka nevejde na slide, zmenši
% ji buď globálně tady, nebo jen v jedné prezentaci stejným řádkem za
% \input{...preambule}:
% \renewcommand{\SPSCodeSize}{\scriptsize}

% --- zkratky pro vkládání kódu ---------------------------------------------
\newcommand{\kod}[1]{\texttt{\small #1}}          % kód uvnitř věty
\newcommand{\kodSoubor}[1]{\lstinputlisting{#1}}  % celá skica ze souboru
\newcommand{\kodSouborRadky}[3]{\lstinputlisting[firstline=#2,lastline=#3]{#1}}

% --- upozornění na slidu ---------------------------------------------------
\newcommand{\pozor}[1]{\textcolor{akcent}{\textbf{#1}}}

% --- metadata předmětu -----------------------------------------------------
\author{Jaroslav Bušek}
\institute{Střední průmyslová škola na Proseku}
% \date{} zůstává prázdné záměrně -- datum a čas posledního překladu se
% zobrazuje automaticky v patě každého slidu (\SPSVersionStamp), takže je
% při promítání hned vidět, jak stará je verze PDF.
\date{}

% ---------------------------------------------------------------------------

\usepackage{iftex}

% --- písma a čeština -------------------------------------------------------
\ifXeTeX
  \usepackage{fontspec}
  \usepackage{polyglossia}
  \setmainlanguage{czech}
  \setmainfont{Latin Modern Roman}
  \setsansfont{Latin Modern Sans}
  \setmonofont{Latin Modern Mono}[Scale=MatchLowercase]
\else
  % Záložní varianta, kdyby někdo přeložil pdfLaTeXem.
  \usepackage[T1]{fontenc}
  \usepackage[utf8]{inputenc}
  \usepackage{lmodern}
  \usepackage[czech]{babel}
\fi

\usepackage{graphicx}
\usepackage{booktabs}
\usepackage{xcolor}
\usepackage{tikz}
\usetikzlibrary{arrows.meta, positioning, shapes.geometric}

% --- motiv SPŠ na Proseku --------------------------------------------------
% Cesta ke složce se stylem a logy. Počítá s tím, že prezentace leží
% v temata/<tema>/prezentace/; zkusí se ale i mělčí zanoření, aby šlo
% přeložit i rozpracovaný soubor jinde v repozitáři.
\IfFileExists{../../../spolecne/styl/beamerthemeSPSProsek.sty}
  {\newcommand{\SPSStylPath}{../../../spolecne/styl}}
  {\IfFileExists{../../spolecne/styl/beamerthemeSPSProsek.sty}
    {\newcommand{\SPSStylPath}{../../spolecne/styl}}
    {\IfFileExists{../spolecne/styl/beamerthemeSPSProsek.sty}
      {\newcommand{\SPSStylPath}{../spolecne/styl}}
      {\newcommand{\SPSStylPath}{.}}}}
\usepackage{\SPSStylPath/beamerthemeSPSProsek}

\setbeamertemplate{footline}[SPSProsek]
\beamertemplatenavigationsymbolsempty

% --- velikost písma ve výpisech kódu ---------------------------------------
% Výchozí je \footnotesize. Když se dlouhá ukázka nevejde na slide, zmenši
% ji buď globálně tady, nebo jen v jedné prezentaci stejným řádkem za
% % ---------------------------------------------------------------------------
% Společná preambule pro prezentace předmětu Programování (3. ročník)
%
% Vzhled řeší motiv beamerthemeSPSProsek.sty ve stejné složce (převzatý
% z šablony předmětů ZPG/SNT na FS ČVUT a přeznačkovaný na SPŠ na Proseku).
% Tady se jen načte a doplní se metadata a makra předmětu.
%
% Překlad: XeLaTeX (kvůli fontspec a českým znakům)
%   latexmk -xelatex prezentace.tex
% nebo jednoduše  .\nastroje\build_prezentace.ps1
%
% Vkládá se z prezentace v temata/<tema>/prezentace/ takto:
%   \input{../../../spolecne/styl/preambule}
% ---------------------------------------------------------------------------

\usepackage{iftex}

% --- písma a čeština -------------------------------------------------------
\ifXeTeX
  \usepackage{fontspec}
  \usepackage{polyglossia}
  \setmainlanguage{czech}
  \setmainfont{Latin Modern Roman}
  \setsansfont{Latin Modern Sans}
  \setmonofont{Latin Modern Mono}[Scale=MatchLowercase]
\else
  % Záložní varianta, kdyby někdo přeložil pdfLaTeXem.
  \usepackage[T1]{fontenc}
  \usepackage[utf8]{inputenc}
  \usepackage{lmodern}
  \usepackage[czech]{babel}
\fi

\usepackage{graphicx}
\usepackage{booktabs}
\usepackage{xcolor}
\usepackage{tikz}
\usetikzlibrary{arrows.meta, positioning, shapes.geometric}

% --- motiv SPŠ na Proseku --------------------------------------------------
% Cesta ke složce se stylem a logy. Počítá s tím, že prezentace leží
% v temata/<tema>/prezentace/; zkusí se ale i mělčí zanoření, aby šlo
% přeložit i rozpracovaný soubor jinde v repozitáři.
\IfFileExists{../../../spolecne/styl/beamerthemeSPSProsek.sty}
  {\newcommand{\SPSStylPath}{../../../spolecne/styl}}
  {\IfFileExists{../../spolecne/styl/beamerthemeSPSProsek.sty}
    {\newcommand{\SPSStylPath}{../../spolecne/styl}}
    {\IfFileExists{../spolecne/styl/beamerthemeSPSProsek.sty}
      {\newcommand{\SPSStylPath}{../spolecne/styl}}
      {\newcommand{\SPSStylPath}{.}}}}
\usepackage{\SPSStylPath/beamerthemeSPSProsek}

\setbeamertemplate{footline}[SPSProsek]
\beamertemplatenavigationsymbolsempty

% --- velikost písma ve výpisech kódu ---------------------------------------
% Výchozí je \footnotesize. Když se dlouhá ukázka nevejde na slide, zmenši
% ji buď globálně tady, nebo jen v jedné prezentaci stejným řádkem za
% \input{...preambule}:
% \renewcommand{\SPSCodeSize}{\scriptsize}

% --- zkratky pro vkládání kódu ---------------------------------------------
\newcommand{\kod}[1]{\texttt{\small #1}}          % kód uvnitř věty
\newcommand{\kodSoubor}[1]{\lstinputlisting{#1}}  % celá skica ze souboru
\newcommand{\kodSouborRadky}[3]{\lstinputlisting[firstline=#2,lastline=#3]{#1}}

% --- upozornění na slidu ---------------------------------------------------
\newcommand{\pozor}[1]{\textcolor{akcent}{\textbf{#1}}}

% --- metadata předmětu -----------------------------------------------------
\author{Jaroslav Bušek}
\institute{Střední průmyslová škola na Proseku}
% \date{} zůstává prázdné záměrně -- datum a čas posledního překladu se
% zobrazuje automaticky v patě každého slidu (\SPSVersionStamp), takže je
% při promítání hned vidět, jak stará je verze PDF.
\date{}
:
% \renewcommand{\SPSCodeSize}{\scriptsize}

% --- zkratky pro vkládání kódu ---------------------------------------------
\newcommand{\kod}[1]{\texttt{\small #1}}          % kód uvnitř věty
\newcommand{\kodSoubor}[1]{\lstinputlisting{#1}}  % celá skica ze souboru
\newcommand{\kodSouborRadky}[3]{\lstinputlisting[firstline=#2,lastline=#3]{#1}}

% --- upozornění na slidu ---------------------------------------------------
\newcommand{\pozor}[1]{\textcolor{akcent}{\textbf{#1}}}

% --- metadata předmětu -----------------------------------------------------
\author{Jaroslav Bušek}
\institute{Střední průmyslová škola na Proseku}
% \date{} zůstává prázdné záměrně -- datum a čas posledního překladu se
% zobrazuje automaticky v patě každého slidu (\SPSVersionStamp), takže je
% při promítání hned vidět, jak stará je verze PDF.
\date{}
:
% \renewcommand{\SPSCodeSize}{\scriptsize}

% --- zkratky pro vkládání kódu ---------------------------------------------
\newcommand{\kod}[1]{\texttt{\small #1}}          % kód uvnitř věty
\newcommand{\kodSoubor}[1]{\lstinputlisting{#1}}  % celá skica ze souboru
\newcommand{\kodSouborRadky}[3]{\lstinputlisting[firstline=#2,lastline=#3]{#1}}

% --- upozornění na slidu ---------------------------------------------------
\newcommand{\pozor}[1]{\textcolor{akcent}{\textbf{#1}}}

% --- metadata předmětu -----------------------------------------------------
\author{Jaroslav Bušek}
\institute{Střední průmyslová škola na Proseku}
% \date{} zůstává prázdné záměrně -- datum a čas posledního překladu se
% zobrazuje automaticky v patě každého slidu (\SPSVersionStamp), takže je
% při promítání hned vidět, jak stará je verze PDF.
\date{}


\renewcommand{\SPSCodeSize}{\scriptsize}
\renewcommand{\headlineA}{Téma 01 -- HW platformy (U1, U2)}

% Strukturu hodiny nesou klíčové snímky (keyframe), ne dělicí snímky sekcí.
\SPSSectionSlidesOff

% Správná odpověď v nabídce. Na prvním překrytí vypadá jako ostatní možnosti,
% červeně se zvýrazní teprve s odkrytím odpovědi (\pause). V režimu handout
% se překrytí slučují, takže v tištěném klíči je zvýrazněná rovnou.
\newcommand{\spravne}[1]{\alt<1>{#1}{\pozor{#1}}}

\title{Hardwarové platformy}
\subtitle{Na jakém stroji náš program běží -- architektury, AVR a ARM \\
  Programování -- 3.\,ročník (Mechatronika), téma 01}

\begin{document}

\maketitle

% ===== KLÍČOVÝ SNÍMEK: cíle hodiny =========================================
\begin{keyframe}{Co dnes zvládneme}
  \begin{itemize}
    \item Vysvětlím, co je \pozor{hardwarová platforma} a proč se stejný program
          nepřeloží pro každý procesor.
    \item Popíšu cestu od zdrojového kódu k programu v paměti mikrokontroleru.
    \item Porovnám \pozor{AVR a ARM} a u zadané úlohy zdůvodním, kterou platformu bych vybral.
  \end{itemize}
  \vfill
  \pozor{Učivo dle ŠVP:} U1 -- možnosti tvorby programů pro různé hardwarové platformy,
  U2 -- vybrané platformy (ARM, AVR), jejich přednosti a omezení.
\end{keyframe}

% ---------------------------------------------------------------------------
\begin{frame}{Kolik mikrokontrolerů jste dnes už použili?}
  \vspace{-2mm}
  \begin{center}
    \obrazek[0.82\textwidth]{../obrazky/01_kde_jsou_mcu}{Koláž: pračka, myš, klíček
      od auta, elektrokolo, nabíječka, ovladač televize -- kde všude je
      mikrokontroler}
  \end{center}
\end{frame}

% ---------------------------------------------------------------------------
\begin{frame}{Proč se tím zabýváme}
  Ročně se vyrobí přes \pozor{25 miliard} mikrokontrolerů -- víc než procesorů
  do PC a mobilů dohromady.

  \vspace{4mm}
  \begin{block}{Co mají všechny společné}
    Jedna úloha, žádný operační systém, cena v jednotkách korun, spotřeba
    v miliampérech, a musí fungovat roky bez restartu.
  \end{block}

  \vspace{4mm}
  \pozor{Náš rok:} naučit se pro takový čip napsat program tak, aby fungoval
  spolehlivě.
\end{frame}

% ===========================================================================
\section{U1 -- Možnosti tvorby programů pro různé platformy}

% ===== KLÍČOVÝ SNÍMEK: cíle části U1 =======================================
\begin{keyframe}{Část 1 (U1): co se dozvíte}
  \small
  \begin{enumerate}\itemsep2pt
    \item Co je hardwarová platforma a co všechno určuje.
    \item Proč procesor nerozumí jazyku C -- a co s tím dělá překladač.
    \item Jak z naší skici vznikne program ve Flash paměti čipu.
    \item Na jakých úrovních se dá pro čip programovat a co která stojí.
    \item Jak jsou v ATmega328P rozdělené paměti a jak se do nich program dostane.
  \end{enumerate}
  \vfill
  \pozor{Nejde o detaily.} Cílem je orientace: vědět, co se kde děje a jaká slova
  znamenají co. Podrobnosti přijdou u konkrétních periferií.
\end{keyframe}

% ---------------------------------------------------------------------------
\begin{frame}{Co je hardwarová platforma}
  \begin{block}{Definice}
    \pozor{Hardwarová platforma} je kombinace procesorového jádra, jeho periferií
    a desky, na které to všechno je. Pro každou platformu vzniká \textbf{jiný}
    výsledný program -- i když zdrojový kód je stejný.
  \end{block}
  \vspace{2mm}
  \small
  \begin{tabular}{@{}p{3.1cm}p{3.2cm}p{4.2cm}@{}}
    \toprule
    Vrstva & Co určuje & U Arduina UNO \\
    \midrule
    \textbf{Architektura jádra} & instrukční sadu, šířku slova, registry
      & AVR, 8 bitů \\
    \textbf{Konkrétní čip} & kolik paměti a jaké periferie
      & ATmega328P \\
    \textbf{Deska} & napájení, konektory, převodník USB
      & Arduino UNO R3 \\
    \bottomrule
  \end{tabular}
\end{frame}

% ---------------------------------------------------------------------------
\begin{frame}{Když přejdu na jinou platformu}
  \begin{columns}[T]
    \begin{column}{0.47\textwidth}
      \begin{exampleblock}{Co zůstane}
        \small
        \begin{itemize}
          \item algoritmus a návrh programu
          \item jazyk C/C++
          \item struktura \kod{setup()} / \kod{loop()}
          \item způsob myšlení o časování a stavech
        \end{itemize}
      \end{exampleblock}
    \end{column}
    \begin{column}{0.47\textwidth}
      \begin{alertblock}{Co se změní}
        \small
        \begin{itemize}
          \item překladač (\kod{avr-gcc} $\to$ \kod{arm-gcc})
          \item názvy registrů a čísla pinů
          \item knihovny pro periferie
          \item napájecí napětí (5 V $\to$ 3,3 V)
        \end{itemize}
      \end{alertblock}
    \end{column}
  \end{columns}
  \vspace{2mm}
  \small
  \pozor{Proto se to učí takhle:} letos na jedné platformě do hloubky. Přechod na
  druhou pak není nové programování, jen nové jméno pro tu samou věc.
\end{frame}

% ---------------------------------------------------------------------------
\begin{frame}[fragile]{Stejný řádek v C, jiné instrukce}
  \begin{center}
    \small \kod{y = y + 1;} \quad kde \kod{y} je bajt v paměti
  \end{center}
  \vspace{-1mm}
  \begin{columns}[T]
    \begin{column}{0.47\textwidth}
      \centering \footnotesize \textbf{AVR} (ATmega328P, 8 bitů)
\begin{vystup}
lds  r24, 0x0100   ; nacti y
subi r24, 0xFF     ; pricti 1
sts  0x0100, r24   ; uloz y
\end{vystup}
    \end{column}
    \begin{column}{0.47\textwidth}
      \centering \footnotesize \textbf{ARM} Cortex-M (32 bitů)
\begin{vystup}
ldr  r3, [pc, #16] ; adresa y
ldrb r2, [r3, #0]  ; nacti y
adds r2, r2, #1    ; pricti 1
strb r2, [r3, #0]  ; uloz y
\end{vystup}
    \end{column}
  \end{columns}
  \vspace{2mm}
  \footnotesize
  Jiné názvy instrukcí, jiný počet registrů, jiná šířka slova. Procesor
  nerozumí jazyku C -- rozumí \pozor{jen své instrukční sadě}. Překladač je ten,
  kdo náš kód do ní přeloží, a pro každou architekturu je to \textbf{jiný překladač}.
  (Výpisy jsou zjednodušené.)
\end{frame}

% ---------------------------------------------------------------------------
\begin{frame}{Křížový překlad -- překládám na PC, běží to na čipu}
  \begin{columns}[T]
    \begin{column}{0.52\textwidth}
      \small
      Na PC nemáme AVR procesor, a přesto pro něj umíme vyrobit program.
      Tomu se říká \pozor{křížový překlad} (cross-compile).

      \vspace{3mm}
      \footnotesize
      \begin{tabular}{@{}ll@{}}
        \toprule
        Cíl & Sada nástrojů \\
        \midrule
        PC (Windows) & \kod{gcc}, \kod{cl.exe} \\
        AVR & \kod{avr-gcc}, \kod{avrdude} \\
        ARM Cortex-M & \kod{arm-none-eabi-gcc} \\
        \bottomrule
      \end{tabular}
    \end{column}
    \begin{column}{0.44\textwidth}
      \small
      Arduino IDE tyto nástroje obsahuje a volá je za nás. Proto se instalace
      „jednoho programu`` postará o celý řetězec.

      \vspace{2mm}
      \pozor{Slovníček:} sada nástrojů pro jednu platformu se souhrnně nazývá
      \textbf{toolchain}.
    \end{column}
  \end{columns}
\end{frame}

% ---------------------------------------------------------------------------
\begin{frame}[fragile]{Kvíz -- proč nejde \kod{program.exe} nahrát do Arduina?}
  \footnotesize
  \begin{itemize}\itemsep2pt
    \item[A)] Jde to, jen se musí přejmenovat na \kod{.hex}.
    \item[B)] Nejde, protože \kod{.exe} je moc velký soubor.
    \item[C)] \spravne{Nejde, protože obsahuje instrukce x86, volání operačního
      systému a počítá s gigabajty RAM.}
    \item[D)] Nejde, protože Arduino umí jen jazyk C, ne C++.
  \end{itemize}

  \pause
  \vspace{2mm}
  \begin{exampleblock}{C)}
    \footnotesize Nic z toho na ATmega328P není: jiná instrukční sada, žádný
    operační systém, 2 kB RAM. Výsledný program je proto vždycky svázaný
    s platformou, pro kterou byl přeložený.
  \end{exampleblock}
\end{frame}

% ---------------------------------------------------------------------------
\begin{frame}{Cesta od naší skici do paměti čipu}
  \vspace{-1mm}
  \begin{center}
    \begin{tikzpicture}[
      font=\scriptsize,
      box/.style={draw=SPSDark, fill=SPSLight, rounded corners=2pt,
                  minimum height=8mm, text width=18mm, align=center, inner sep=0.8mm},
      fil/.style={draw=SPSRed, thick, fill=SPSRed!8, rounded corners=2pt,
                  minimum height=7mm, text width=14mm, align=center, inner sep=0.8mm},
      arr/.style={-Stealth, thick, SPSGrey}]

      \node[fil] (ino)  at (0,0)   {\texttt{skica.ino}};
      \node[box] (pre)  at (2.5,0) {preprocesor\\[-0.6mm]{\tiny \texttt{\#include}}};
      \node[box] (kom)  at (5.1,0) {překladač\\[-0.6mm]{\tiny \texttt{avr-g++}}};
      \node[fil] (obj)  at (7.5,0) {\texttt{.o}};

      \node[box] (lnk)  at (0,-1.9)   {linker\\[-0.6mm]{\tiny + knihovny}};
      \node[fil] (elf)  at (2.5,-1.9) {\texttt{.elf}};
      \node[box] (hex)  at (5.1,-1.9) {\texttt{avr-objcopy}};
      \node[fil] (fhex) at (7.5,-1.9) {\texttt{.hex}};
      \node[box, fill=SPSRed!15, draw=SPSRed, thick, text width=20mm]
            (fl) at (10.1,-1.9) {\texttt{avrdude}\\[-0.6mm]{\tiny do Flash}};

      \draw[arr] (ino) -- (pre);
      \draw[arr] (pre) -- (kom);
      \draw[arr] (kom) -- (obj);
      \draw[arr] (obj.south) -- ++(0,-0.5) -| (lnk.north);
      \draw[arr] (lnk) -- (elf);
      \draw[arr] (elf) -- (hex);
      \draw[arr] (hex) -- (fhex);
      \draw[arr] (fhex) -- (fl);
    \end{tikzpicture}
  \end{center}
  \vspace{2mm}
  \footnotesize
  V Arduino IDE se tohle všechno schová za jedno tlačítko. Když se objeví chyba,
  je ale dobré vědět, \pozor{ve kterém kroku} vznikla: chyba v syntaxi je
  z překladače, „undefined reference`` z linkeru a „nelze otevřít port``
  z nahrávání.
\end{frame}

% ---------------------------------------------------------------------------
\begin{frame}[fragile]{Co Arduino IDE přidá k našemu kódu}
  \begin{columns}[T]
    \begin{column}{0.5\textwidth}
\begin{lstlisting}[style=arduino_bez_cisel]
// tohle do skici nepiseme,
// pridá se to samo:
#include <Arduino.h>

void blikni();   // prototypy

int main(void) {
  init();
  setup();
  for (;;) { loop(); }
}
\end{lstlisting}
    \end{column}
    \begin{column}{0.46\textwidth}
      \small
      Proto v Arduinu funguje volání funkce, která je v souboru \textbf{níž} --
      prototyp za nás doplní IDE.

      \vspace{2mm}
      \pozor{Důsledek:} skica \kod{.ino} \textbf{není čisté C}. Je to C++ s několika
      doplňky, a běžným \kod{gcc} se nepřeloží.

      \vspace{2mm}
      \scriptsize
      V bloku 15 proto budeme kód psát do \kod{.h} a \kod{.cpp} -- tam už žádná
      magie není a platí normální pravidla jazyka.
    \end{column}
  \end{columns}
\end{frame}

% ---------------------------------------------------------------------------
\begin{frame}{Možnosti tvorby programů (U1)}
  \small Pro stejný čip lze psát na několika úrovních. Volba je vždy kompromis:
  \vspace{1mm}

  \scriptsize
  \begin{tabular}{@{}p{2.4cm}p{3.7cm}p{1.6cm}p{1.6cm}@{}}
    \toprule
    Úroveň & Jak to vypadá & Rychlost & Pracnost \\
    \midrule
    \textbf{Asembler} & jednotlivé instrukce & nejlepší & velmi vysoká \\
    \textbf{C/C++ s registry} & \kod{PORTB |= (1<<5);} & výborná & vysoká \\
    \textbf{C/C++ s HAL} & \kod{HAL\_GPIO\_WritePin(...)} & dobrá & střední \\
    \textbf{Arduino} & \kod{digitalWrite(13, HIGH);} & horší & nízká \\
    \textbf{MicroPython} & \kod{pin.value(1)} & nejhorší & nejnižší \\
    \bottomrule
  \end{tabular}
  \vspace{2mm}

  \footnotesize
  \pozor{Letos:} Arduino framework -- abychom se soustředili na návrh programu.
  V bloku 03 se podíváme i „pod něj``, na registry: \kod{digitalWrite()} je
  desítky instrukcí tam, kde by stačila jedna.
\end{frame}

% ---------------------------------------------------------------------------
\begin{frame}[fragile]{Co je HAL -- vrstva abstrakce hardwaru}
  \begin{block}{\kod{HAL} = Hardware Abstraction Layer, vrstva abstrakce hardwaru}
    \footnotesize
    Knihovna, obvykle od \textbf{výrobce čipu}, která schová práci s registry za
    funkce s čitelnými názvy. Program pak neříká \emph{„nastav pátý bit v registru
    PORTB``}, ale \emph{„nastav pin PB5 do jedničky``}.
  \end{block}
  \begin{columns}[T]
    \begin{column}{0.52\textwidth}
\begin{lstlisting}[style=arduino_bez_cisel]
// primo registr (bez HAL)
PORTB |= (1 << 5);

// pres HAL (STM32)
HAL_GPIO_WritePin(GPIOB,
                  GPIO_PIN_5,
                  GPIO_PIN_SET);
\end{lstlisting}
    \end{column}
    \begin{column}{0.44\textwidth}
      \footnotesize
      \textbf{Co za to dostanu:} kód se dá přenést na jiný čip téže řady bez
      přepisování a je čitelnější.

      \vspace{2mm}
      \textbf{Co to stojí:} funkce navíc kontroluje parametry, takže je
      \pozor{pomalejší a delší} než jeden řádek s registrem.
    \end{column}
  \end{columns}
  \scriptsize
  \pozor{Souvislost:} \kod{digitalWrite()} je taky HAL, jen o úroveň výš.
\end{frame}

% ---------------------------------------------------------------------------
\begin{frame}{Dvě architektury pamětí}
  \begin{columns}[T]
    \begin{column}{0.48\textwidth}
      \centering \footnotesize \textbf{Harvardská} -- AVR
      \vspace{1mm}

      \begin{tikzpicture}[font=\tiny,
        m/.style={draw=SPSDark, fill=SPSLight, rounded corners=2pt,
                  minimum width=20mm, minimum height=7mm, align=center},
        c/.style={draw=SPSRed, thick, fill=SPSRed!10, rounded corners=2pt,
                  minimum width=16mm, minimum height=8mm},
        a/.style={-Stealth, thick, SPSGrey}]
        \node[m] (f) at (0,1.2) {Flash\\program};
        \node[c] (cpu) at (0,0) {jádro};
        \node[m] (s) at (0,-1.2) {SRAM\\data};
        \draw[a] (cpu) -- (f) node[midway,right,SPSGrey] {\,sběrnice 1};
        \draw[a] (cpu) -- (s) node[midway,right,SPSGrey] {\,sběrnice 2};
      \end{tikzpicture}

      \vspace{1mm}
      \scriptsize Oddělené sběrnice \textbf{i oddělené adresy}. Adresa 0x0100
      ve Flash a v SRAM jsou dvě různá místa.
      \vspace{1mm}
    \end{column}
    \begin{column}{0.48\textwidth}
      \centering \footnotesize \textbf{Jednotný prostor} -- ARM Cortex-M
      \vspace{1mm}

      \begin{tikzpicture}[font=\tiny,
        m/.style={draw=SPSDark, fill=SPSLight, rounded corners=2pt,
                  minimum width=30mm, minimum height=7mm, align=center},
        c/.style={draw=SPSRed, thick, fill=SPSRed!10, rounded corners=2pt,
                  minimum width=16mm, minimum height=8mm},
        a/.style={-Stealth, thick, SPSGrey}]
        \node[m] (f) at (0,1.2) {jedna adresní mapa\\Flash + SRAM + periferie};
        \node[c] (cpu) at (0,0) {jádro};
        \node[m, minimum width=30mm] (s) at (0,-1.2) {ukazatel může mířit\\kamkoliv};
        \draw[a] (cpu) -- (f);
        \draw[a] (cpu) -- (s);
      \end{tikzpicture}

      \vspace{1mm}
      \scriptsize Jeden adresní prostor. Ukazatel na konstantu ve Flash se
      používá stejně jako na proměnnou. (Sběrnice uvnitř jádra oddělené být
      mohou i nemusí -- podrobněji v části U2.)
    \end{column}
  \end{columns}
  \vspace{1mm}
  \scriptsize
  \pozor{Důsledek pro nás:} na AVR se textové konstanty kopírují při startu do SRAM.
  Proto \kod{Serial.println(F("text"))} -- \kod{F()} je nechá ve Flash (blok 09).
\end{frame}

% ---------------------------------------------------------------------------
\begin{frame}{Paměti ATmega328P (U7)}
  \vspace{2mm}
  \begin{center}
    \begin{tikzpicture}[font=\scriptsize]
      % 1 kB = 0.25 cm  ->  Flash 32 kB = 8 cm, SRAM 2 kB = 0.5 cm, EEPROM 1 kB = 0.25 cm
      \draw[thick, draw=SPSDark, fill=SPSRed!15] (0,0)     rectangle ++(8,0.55);
      \node[anchor=west] at (8.2,0.275) {\textbf{Flash 32 kB}};
      \draw[thick, draw=SPSDark, fill=SPSLight]  (0,-0.85) rectangle ++(0.5,0.55);
      \node[anchor=west] at (0.7,-0.575) {\textbf{SRAM 2 kB}};
      \draw[thick, draw=SPSDark, fill=SPSGrey!25] (0,-1.7) rectangle ++(0.25,0.55);
      \node[anchor=west] at (0.45,-1.425) {\textbf{EEPROM 1 kB}};
    \end{tikzpicture}
  \end{center}
  \vspace{3mm}
  \small
  \pozor{To je celá paměť, se kterou letos pracujeme.} Program se vejde téměř vždy;
  co dělá potíže, je \textbf{SRAM} -- těch 2048 bajtů se spotřebuje překvapivě rychle.
  Proto se od začátku hlídá, kolik toho který program zabere.
\end{frame}

% ---------------------------------------------------------------------------
\begin{frame}{Co je v které paměti}
  \footnotesize
  \begin{tabular}{@{}p{1.9cm}p{4.2cm}p{4.4cm}@{}}
    \toprule
    Paměť & Co v ní je & Vlastnosti \\
    \midrule
    \textbf{Flash} & náš program, textové konstanty, 0,5 kB bootloader
      & po vypnutí zůstane, přepisuje se jen při nahrávání \\
    \addlinespace
    \textbf{SRAM} & proměnné, pole, zásobník, buffery \kod{Serial}
      & po vypnutí se \pozor{smaže}, rychlá, a je jí málo \\
    \addlinespace
    \textbf{EEPROM} & nastavení, které má přežít vypnutí
      & zůstane, zapisuje se pomalu a jen \textasciitilde 100\,000 krát \\
    \bottomrule
  \end{tabular}
  \vspace{3mm}

  \small
  \pozor{Typická otázka u zkoušení:} kam patří text \kod{"Teplota: "}? Do Flash --
  ale bez \kod{F()} se při startu zkopíruje i do SRAM, takže zabere obojí.
\end{frame}

% ---------------------------------------------------------------------------
\begin{frame}{Jak se program dostane do Flash (U6)}
  \vspace{1mm}
  \begin{center}
    \begin{tikzpicture}[font=\scriptsize,
      b/.style={draw=SPSDark, fill=SPSLight, rounded corners=2pt,
                minimum height=7.5mm, text width=20mm, align=center, inner sep=0.8mm},
      r/.style={draw=SPSRed, thick, fill=SPSRed!10, rounded corners=2pt,
                minimum height=7.5mm, text width=20mm, align=center, inner sep=0.8mm},
      a/.style={-Stealth, thick, SPSGrey}]

      \node[font=\scriptsize\bfseries, anchor=west] at (-1.0,0.75)
        {1. Bootloader -- běžné nahrávání};
      \node[b] (pc1) at (0,0)   {PC, USB kabel};
      \node[b] (cv)  at (3.0,0) {převodník\\{\tiny USB\,$\to$\,sériová linka}};
      \node[r] (bl)  at (6.0,0) {bootloader\\v čipu};
      \node[b] (fl1) at (9.0,0) {Flash};
      \draw[a] (pc1) -- (cv); \draw[a] (cv) -- (bl); \draw[a] (bl) -- (fl1);

      \node[font=\scriptsize\bfseries, anchor=west] at (-1.0,-1.3)
        {2. ICSP -- programátorem přímo};
      \node[b] (pc2) at (0,-2.05)   {PC};
      \node[b] (pg)  at (3.0,-2.05) {programátor\\{\tiny např. USBasp}};
      \node[r] (icsp)at (6.0,-2.05) {piny ICSP\\{\tiny SPI + RESET}};
      \node[b] (fl2) at (9.0,-2.05) {Flash};
      \draw[a] (pc2) -- (pg); \draw[a] (pg) -- (icsp); \draw[a] (icsp) -- (fl2);
    \end{tikzpicture}
  \end{center}
  \vspace{2mm}
  \small
  Bootloader je \pozor{program, který nahrává jiný program}. Po resetu chvíli čeká,
  jestli po sériové lince nepřijde nový kód; když nepřijde, spustí ten starý.
\end{frame}

% ---------------------------------------------------------------------------
\begin{frame}{Bootloader nebo ICSP -- kdy co}
  \footnotesize
  \begin{tabular}{@{}p{2.2cm}p{4.0cm}p{4.2cm}@{}}
    \toprule
     & \textbf{Bootloader} & \textbf{ICSP} \\
    \midrule
    Potřebuje    & jen USB kabel & programátor (USBasp, druhé Arduino) \\
    \addlinespace
    Cena za to   & ubere 0,5 kB Flash, zdrží start & nic \\
    \addlinespace
    Nový, prázdný čip & nenahraje (bootloader tam ještě není) & nahraje \\
    \addlinespace
    Kdy se použije & běžná výuka a vývoj & výroba, oživení čipu, potřeba celé Flash \\
    \bottomrule
  \end{tabular}
  \vspace{3mm}

  \small
  \pozor{Praktický dopad na nás:} celý rok nahráváme přes bootloader a USB.
  ICSP stačí znát -- a vědět, že těch 0,5 kB Flash chybí, když se program
  „skoro`` vejde.
\end{frame}

% ===== KLÍČOVÝ SNÍMEK: shrnutí části U1 ====================================
\begin{keyframe}{Shrnutí části 1 (U1)}
  \footnotesize
  \begin{itemize}\itemsep2pt
    \item \pozor{Platforma} $=$ jádro $+$ čip $+$ deska. Každé jádro má vlastní
          instrukční sadu, a tedy vlastní překladač.
    \item Program se vyrábí \pozor{křížovým překladem} na PC:
          \kod{.ino} $\to$ preprocesor $\to$ překlad $\to$ linker $\to$ \kod{.hex}
          $\to$ nahrání do Flash.
    \item Skica \kod{.ino} není čisté C -- IDE k ní přidá \kod{Arduino.h},
          prototypy a \kod{main()}.
    \item Programovat se dá od \pozor{asembleru} po \pozor{MicroPython}; Arduino
          framework je kompromis mezi pracností a rychlostí.
    \item AVR je \pozor{harvardská} architektura -- Flash a SRAM mají oddělené adresy.
    \item Flash 32 kB zůstane, \pozor{SRAM 2 kB} se maže a je jí málo,
          EEPROM 1 kB je pomalá.
    \item Nahrává se přes \pozor{bootloader} (jen USB, ubere 0,5 kB) nebo
          \pozor{ICSP} (programátor, uvolní celou Flash).
  \end{itemize}
\end{keyframe}

% ===========================================================================
\section{U2 -- Vybrané platformy: AVR a ARM}

% ===== KLÍČOVÝ SNÍMEK: cíle části U2 =======================================
\begin{keyframe}{Část 2 (U2): co se dozvíte}
  \small
  \begin{enumerate}\itemsep2pt
    \item Co je AVR a co je ARM Cortex-M.
    \item Jaké má každá z nich přednosti a jaká omezení.
    \item Čím se liší uvnitř -- sada instrukcí, paměť, registry.
    \item Podle čeho se platforma pro konkrétní úlohu vybírá.
    \item Co se v téhle oblasti běžně plete -- a proč na tom záleží.
  \end{enumerate}
  \vfill
  \pozor{Na konci této části} byste měli u zadané úlohy umět říct „vzal bych tohle,
  protože\dots`` a uvést dva konkrétní důvody.
\end{keyframe}

% ---------------------------------------------------------------------------
\begin{frame}{Desky, o kterých bude řeč}
  \vspace{-2mm}
  \begin{center}
    \obrazek[0.82\textwidth]{../obrazky/02_rodina_desek}{Fotografie desek vedle sebe:
      Arduino UNO, Arduino Nano, Raspberry Pi Pico, STM32 „Blue Pill``, ESP32}
  \end{center}
\end{frame}

% ---------------------------------------------------------------------------
\begin{frame}{AVR -- na čem se letos učíme}
  \begin{columns}[T]
    \begin{column}{0.52\textwidth}
      \small
      8bitová architektura firmy Atmel, na trhu \textbf{od roku 1996}.
      Dnes ji vlastní Microchip a pořád se vyrábí.

      \vspace{3mm}
      Právě ta dlouhá životnost je důvod, proč je kolem ní tolik hotových
      knihoven, návodů a odpovědí na fóru -- a proč se na ní dobře začíná.
    \end{column}
    \begin{column}{0.44\textwidth}
      \footnotesize
      \textbf{ATmega328P v Arduinu UNO}
      \begin{itemize}
        \item 8 bitů, 16 MHz
        \item 32 kB Flash, 2 kB SRAM, 1 kB EEPROM
        \item 20 I/O pinů, ADC 10 bitů
        \item 3 časovače
        \item napájení 1,8--5,5 V
      \end{itemize}
    \end{column}
  \end{columns}
\end{frame}

% ---------------------------------------------------------------------------
\begin{frame}{AVR -- přednosti a omezení}
  \begin{columns}[T]
    \begin{column}{0.48\textwidth}
      \begin{exampleblock}{Přednosti}
        \footnotesize
        \begin{itemize}
          \item snese 5 V i kolísavé napájení
          \item čitelná dokumentace (\textasciitilde 300 stran)
          \item jedna instrukce za takt -- časování se dá spočítat
          \item v patici DIP, dá se vyměnit
          \item obrovská nabídka knihoven a návodů
        \end{itemize}
      \end{exampleblock}
    \end{column}
    \begin{column}{0.48\textwidth}
      \begin{alertblock}{Omezení}
        \footnotesize
        \begin{itemize}
          \item 8 bitů -- \kod{float} se počítá softwarově a pomalu
          \item 2 kB SRAM
          \item 16 MHz
          \item žádná Wi-Fi ani USB
          \item ADC jen 10 bitů
        \end{itemize}
      \end{alertblock}
    \end{column}
  \end{columns}
  \vspace{1mm}
  \small
  \pozor{Proč se na tom učí:} málo možností znamená, že je vidět, co program dělá.
  Na 8 bitech a 2 kB se nedá „zamést`` neefektivní řešení pod koberec.
\end{frame}

% ---------------------------------------------------------------------------
\begin{frame}{ARM Cortex-M -- kam se jde dál}
  \small
  ARM sám čipy nevyrábí -- \pozor{licencuje jádro} a čipy staví desítky firem
  (ST, NXP, Nordic, Raspberry Pi). Proto je nabídka tak široká.

  \vspace{3mm}
  \footnotesize
  \begin{tabular}{@{}p{3.2cm}p{3.2cm}p{4.4cm}@{}}
    \toprule
    Čip & Jádro & Parametry \\
    \midrule
    STM32F103 („Blue Pill``) & Cortex-M3, 72 MHz & 64 kB Flash, 20 kB SRAM, ADC 12 b \\
    \addlinespace
    RP2040 (Pi Pico) & 2$\times$ Cortex-M0+, 133 MHz & 264 kB SRAM, Flash externí \\
    \addlinespace
    nRF52840 & Cortex-M4F, 64 MHz & Bluetooth, 256 kB SRAM \\
    \bottomrule
  \end{tabular}
\end{frame}

% ---------------------------------------------------------------------------
\begin{frame}{ARM -- přednosti a omezení}
  \begin{columns}[T]
    \begin{column}{0.48\textwidth}
      \begin{exampleblock}{Přednosti}
        \footnotesize
        \begin{itemize}
          \item 32 bitů, desítky až stovky MHz
          \item mnohem víc paměti
          \item bohaté periferie (USB, CAN, Ethernet)
          \item hardwarový \kod{float} u M4F
          \item hluboké režimy spánku
        \end{itemize}
      \end{exampleblock}
    \end{column}
    \begin{column}{0.48\textwidth}
      \begin{alertblock}{Omezení}
        \footnotesize
        \begin{itemize}
          \item jen 3,3 V -- 5 V na pinu čip zničí
          \item složitější rozběh (hodiny, napájecí domény)
          \item datasheet má tisíc stran
          \item čip obvykle nelze vyměnit
        \end{itemize}
      \end{alertblock}
    \end{column}
  \end{columns}
  \vspace{1mm}
  \small
  \pozor{Nejčastější zničený čip ve školní laboratoři:} ARM, na který někdo přivedl
  5 V ze staré desky. Napájecí úroveň se kontroluje \textbf{před} zapojením.
\end{frame}

% ---------------------------------------------------------------------------
\begin{frame}{Obě jsou RISC -- a přesto jiné}
  \small
  \pozor{RISC} (\textit{Reduced Instruction Set Computer}) znamená málo
  jednoduchých instrukcí, zato rychlých. Opak je CISC -- procesory x86 v PC.
  \vspace{2mm}

  \scriptsize
  \begin{tabular}{@{}p{2.1cm}p{4.3cm}p{4.3cm}@{}}
    \toprule
     & \textbf{AVR} & \textbf{ARM} \\
    \midrule
    Sada instrukcí & RISC, 8 bitů; většina instrukcí za 1 takt
      & RISC, 32 bitů; sady ARM, Thumb, Thumb-2 \\
    \addlinespace
    Paměť & striktně harvardská
      & podle jádra: von Neumannova (M0, M0+) nebo modifikovaná harvardská (M3, M4) \\
    \addlinespace
    Registry & 32 registrů po 8 bitech & 16 registrů po 32 bitech \\
    \addlinespace
    Licence & proprietární, Microchip (dříve Atmel)
      & ARM Ltd. jádro jen navrhuje a licencuje (ST, NXP, Apple) \\
    \bottomrule
  \end{tabular}
  \vspace{2mm}

  \footnotesize
  \pozor{Cortex-M umí jen Thumb-2}, ne původní 32bitovou sadu ARM. Instrukce jsou
  kratší, takže se program vejde do menší Flash.
\end{frame}

% ---------------------------------------------------------------------------
\begin{frame}{Registry -- tam, kde se opravdu počítá}
  \begin{block}{Obě architektury jsou typu \textit{load/store}}
    \footnotesize
    Aritmetika neumí sáhnout přímo do RAM. Hodnota se musí nejdřív \textbf{načíst}
    do registru, tam se s ní počítá, a výsledek se \textbf{uloží} zpět.
  \end{block}
  \vspace{1mm}
  \begin{columns}[T]
    \begin{column}{0.47\textwidth}
      \footnotesize \textbf{AVR}
      \begin{itemize}
        \scriptsize
        \item 32 registrů po 8 bitech, \kod{R0}--\kod{R31}
        \item poslední tři páry slouží jako 16bitové ukazatele
              \kod{X}, \kod{Y}, \kod{Z}
        \item registry visí přímo na ALU -- operace nad dvěma registry
              proběhne za jeden takt
      \end{itemize}
    \end{column}
    \begin{column}{0.47\textwidth}
      \footnotesize \textbf{ARM Cortex-M}
      \begin{itemize}
        \scriptsize
        \item 16 registrů po 32 bitech, \kod{R0}--\kod{R15}
        \item obecných je jich ale jen 13 (\kod{R0}--\kod{R12})
        \item \kod{R13} $=$ \kod{SP} (vrchol zásobníku),
              \kod{R14} $=$ \kod{LR} (návratová adresa),
              \kod{R15} $=$ \kod{PC} (další instrukce)
      \end{itemize}
    \end{column}
  \end{columns}
  \vspace{1mm}
  \footnotesize
  \pozor{Zpátky k tomu výpisu na začátku hodiny:} \kod{lds} a \kod{sts} u AVR,
  \kod{ldr} a \kod{str} u ARM -- to jsou přesně ty \textit{load} a \textit{store}.
  Sčítalo se až potom, a jen mezi registry.
\end{frame}

% ---------------------------------------------------------------------------
\begin{frame}{Srovnání -- co rozhoduje při výběru}
  \footnotesize
  \begin{tabular}{@{}p{3.0cm}p{3.6cm}p{4.2cm}@{}}
    \toprule
     & \textbf{AVR} (ATmega328P) & \textbf{ARM} (např. RP2040) \\
    \midrule
    Šířka slova    & 8 bitů                & 32 bitů \\
    Frekvence      & 16 MHz                & 133 MHz \\
    SRAM           & 2 kB                  & 264 kB \\
    Napájení       & 1,8--5,5 V (snese 5 V)& 1,8--3,6 V (jen 3,3 V) \\
    \kod{float}    & softwarově, pomalu    & rychle (u M4F hardwarově) \\
    Periferie      & UART, I2C, SPI, ADC   & + USB, PWM bloky, PIO, DMA \\
    Rozběh čipu    & funguje hned          & nastavit hodiny a napájení \\
    Dokumentace    & \textasciitilde 300 stran, čitelná & \textasciitilde 600--1000 stran \\
    \bottomrule
  \end{tabular}
  \vspace{2mm}

  \small
  \pozor{Pravidlo:} nevybírá se „lepší`` čip, ale \textbf{nejmenší, který úlohu zvládne}.
  Rozhoduje spotřeba, napětí periferií, potřebná paměť a cena v sérii.
\end{frame}

% ---------------------------------------------------------------------------
\begin{frame}{Kvíz -- co byste vybrali a proč?}
  \small
  \begin{enumerate}\itemsep3pt
    \item Blikající světlo na kolo. Dvě LED, tlačítko, běží z knoflíkové baterie
          několik měsíců.
    \item Kamera u vchodu, která posílá snímek do mobilu přes Wi-Fi.
    \item Regulace otáček motoru s enkodérem, 1000 přepočtů za sekundu s \kod{float}.
  \end{enumerate}

  \pause
  \vspace{2mm}
  \footnotesize
  \begin{tabular}{@{}p{0.4cm}p{3.2cm}p{6.8cm}@{}}
    \toprule
    & Volba & Proč \\
    \midrule
    1. & \textbf{AVR} (i menší ATtiny) & úloha je triviální, rozhoduje spotřeba
         ve spánku a cena \\
    2. & \textbf{ESP32} & Wi-Fi musí být na čipu; samotné ARM jádro to neřeší \\
    3. & \textbf{ARM} Cortex-M4F & 1000$\times$ za sekundu \kod{float} 8bitový AVR
         nestihne \\
    \bottomrule
  \end{tabular}
\end{frame}

% ---------------------------------------------------------------------------
\begin{frame}{Pozor na časté záměny}
  \footnotesize
  \begin{tabular}{@{}p{4.4cm}p{6.4cm}@{}}
    \toprule
    Říká se & Ve skutečnosti \\
    \midrule
    „ESP32 je ARM`` & není -- má jádro \textbf{Xtensa} (novější řady i RISC-V) \\
    \addlinespace
    „Arduino je procesor`` & Arduino je \textbf{prostředí a knihovny};
      běží na AVR, ARM i ESP32 \\
    \addlinespace
    „Raspberry Pi a Pico je totéž`` & Pi je \textbf{počítač s Linuxem},
      Pico je \textbf{mikrokontroler} \\
    \addlinespace
    „procesor`` $=$ „mikrokontroler`` & mikrokontroler má jádro,
      paměti \textbf{i periferie} na jednom čipu \\
    \bottomrule
  \end{tabular}
  \vspace{3mm}

  \small
  \pozor{Proč na tom záleží:} podle těchto slov se hledá dokumentace a knihovny.
  Kdo hledá „ARM knihovnu`` pro ESP32, nenajde nic.
\end{frame}

% ---------------------------------------------------------------------------
% REZERVA – když se nestíhá, tento slide vypustit.
\begin{frame}{Čipy zblízka \small (rezerva)}
  \begin{columns}[T]
    \begin{column}{0.48\textwidth}
      \obrazek[\textwidth]{../obrazky/03_cipy_avr_arm}{Fotografie vedle sebe:
        ATmega328P v patici DIP-28 a čip ARM v pouzdře LQFP, se stejným
        měřítkem a počtem vývodů}
    \end{column}
    \begin{column}{0.48\textwidth}
      \obrazek[\textwidth]{../obrazky/04_deska_uno_popis}{Deska Arduino UNO
        s vyznačenými částmi: ATmega328P, převodník USB-UART, krystal 16 MHz,
        stabilizátor 5 V, konektor ICSP, tlačítko RESET, LED na pinu 13}
    \end{column}
  \end{columns}
  \vspace{1mm}
  \footnotesize
  \pozor{Na co se dívat:} kolik má čip vývodů (tolik nejvýš periferií),
  v jakém je pouzdře (DIP se dá vyměnit, LQFP už ne) a co všechno je na desce
  \textbf{kolem} čipu -- napájení, převodník, krystal. Příště si to celé
  proměříme multimetrem.
\end{frame}

% ===== KLÍČOVÝ SNÍMEK: shrnutí hodiny ======================================
\begin{keyframe}{Shrnutí hodiny}
  \footnotesize
  \begin{itemize}\itemsep2pt
    \item \pozor{Platforma} $=$ jádro $+$ čip $+$ deska. Pro každé jádro vzniká jiný
          výsledný program, protože každé má \textbf{jinou instrukční sadu}.
    \item Program vzniká \pozor{křížovým překladem} na PC a nahrává se přes
          \pozor{bootloader} nebo \pozor{ICSP}.
    \item ATmega328P: 32 kB Flash, \pozor{2 kB SRAM}, 1 kB EEPROM -- a to je všechno.
    \item AVR: 8 bitů, 5 V, málo paměti, jednoduché a spolehlivé.
          ARM: 32 bitů, 3,3 V, hodně paměti a periferií, složitější rozběh.
    \item Vybírá se \pozor{nejmenší čip, který úlohu zvládne}.
  \end{itemize}
  \vfill
  \pozor{Příště:} Arduino IDE, první nahraná skica, vlastní repozitář a první
  \kod{commit}.
\end{keyframe}

\end{document}
